\section{Antecedentes del proyecto}

\subsection{Situación previa}
\begin{indentedsubsection}
Perú Flores Travel Agency Tours Operador E.I.R.L. es una agencia con una sólida trayectoria en el sector de turismo receptivo, especializada en organizar paquetes turísticos y servicios personalizados hacia destinos como Cusco, Arequipa y Puno. Con un fuerte enfoque en el mercado extranjero, destacando principalmente la atención a turistas brasileños, la empresa ha logrado sostenerse a lo largo de los años gracias al esfuerzo directo de su gerencia general, liderada por Silveria Flores.

Sin embargo, hasta antes de la intervención tecnológica, la agencia operaba bajo un modelo de gestión tradicional, caracterizado por una digitalización nula o empírica:

\begin{itemize}
    \item \textbf{Captación y comunicación informal:} el primer contacto con los turistas internacionales, así como las consultas y confirmaciones, se realizaban de manera casi exclusiva a través de chats de WhatsApp.
    \item \textbf{Gestión de datos en formatos físicos y aislados:} la base de datos de los pasajeros se gestionaba manualmente, anotando los registros en cuadernos físicos para luego intentar transcribirlos de forma esporádica a la computadora.
    \item \textbf{Armado de paquetes artesanal:} la creación de cotizaciones dependía de buscar plantillas antiguas en Excel (creadas durante la pandemia) o documentos de Word, los cuales se copiaban, pegaban y modificaban manualmente cambiando fechas y nombres.
    \item \textbf{Pagos y finanzas descentralizadas:} al carecer de una pasarela de pagos en línea, la agencia dependía netamente de depósitos internacionales por Western Union o transferencias bancarias directas, requiriendo un control de abonos totalmente manual.
\end{itemize}
\end{indentedsubsection}

\subsection{Problemática identificada}
\begin{indentedsubsection}
El diagnóstico técnico y el relevamiento de procesos revelaron brechas críticas (``silos de información'') que obstaculizaban la escalabilidad del negocio y generaban un desgaste operativo significativo para la gerencia:

\begin{itemize}
    \item \textbf{A. Cuellos de botella y lentitud en cotizaciones:} el proceso de presupuestar un paquete turístico tomaba un tiempo excesivo (entre 45 y 60 minutos por cotización). La gerencia debía calcular manualmente los costos prorrateados (por ejemplo, dividir el costo fijo de 40 soles de transporte y 50 soles de guía entre el número de pasajeros) y sumar los ingresos individuales (como el BTG o Machu Picchu). Esta carga operativa artesanal limitaba la capacidad de respuesta rápida ante nuevos prospectos.
    \item \textbf{B. Fragmentación y vulnerabilidad de la información:} la dependencia de cuadernos físicos y hojas de cálculo dispersas generaba un alto riesgo de pérdida de información y una nula trazabilidad comercial. No existía un historial centralizado que permitiera saber cuántas veces un cliente había viajado con la agencia o recuperar fácilmente cotizaciones pasadas sin tener que buscar manualmente entre múltiples archivos.
    \item \textbf{C. Duplicidad de costos y desconfianza tecnológica:} la empresa mantenía pagos anuales por servicios de dominio y hosting web que no estaban integrados a su flujo operativo y que resultaban ineficientes (enlaces rotos, videos pesados que retrasaban la carga). A esto se sumaban malas experiencias previas con asesorías tecnológicas y consultores que no cumplían con las reuniones ni entregaban resultados tangibles, lo que generaba en la gerencia una justificada frustración y resistencia a incurrir en dobles gastos.
    \item \textbf{D. Ausencia de métricas y retroalimentación (Feedback):} aunque la agencia brindaba un servicio de alta calidad, carecía de una herramienta sistematizada para recopilar las opiniones de los turistas al finalizar el viaje. La gerencia dependía de terceros o procesos manuales para enviar formularios, perdiendo la oportunidad de capitalizar los testimonios positivos para generar un ``boca a boca'' digital que atrajera nuevos clientes.
\end{itemize}

En este contexto, la implementación de Odoo se concibe como una respuesta directa para centralizar la información, erradicar los cálculos manuales y consolidar un ecosistema digital que le devuelva a la gerencia el control total de su agencia sin depender de cuadernos ni herramientas desconectadas.
\end{indentedsubsection}
